% !TEX root = ../aa_SoM_Chapters.tex

% ö  \equal  \%c3\%b6
% ä  \equal  \%C3\%A4

\xdef\sectiontitle{\Isectiontitle} %aiheuttama

%%%%%%%%%%%%%%%
\begin{frame}[label=1_1_a]%[label \equal 1_0_TOC1]
\frametitle{\texorpdfstring{\culine[\sectiontitleulinecolor]{\sectiontitle}}{\sectiontitle} }

\vspace{0.5cm}

\begin{itemize}[leftmargin = 0.5cm,itemsep = 2mm]
    \item[1.1]<1-> \hyperlink{1_1_a}{\Alert[blue]{\IaSubsectiontitle}}
    \item[1.2]<1-> \hyperlink{1_2_a}{\IbSubsectiontitle}	
    \item[1.3]<1-> \hyperlink{1_3_a}{\IcSubsectiontitle}	
    \item[1.4]<1-> \hyperlink{1_4_a}{\IdSubsectiontitle}
    \item[1.5]<1-> \hyperlink{1_5_a}{\IeSubsectiontitle}
    \item[1.6]<1-> \hyperlink{1_6_a}{\IfSubsectiontitle}
    \item[1.7]<1-> \hyperlink{1_7_a}{\IgSubsectiontitle}
\end{itemize}

\vspace{-1cm}
\begin{itemize}[leftmargin = 9.5cm,itemsep = 2mm]
    \item[1.8]<1-> \hyperlink{1_8_a}{\IhSubsectiontitle}
	\item[1.9]<1-> \hyperlink{1_9_a}{\IiSubsectiontitle}
	\item[1.10]<1-> \hyperlink{1_10_a}{\IjSubsectiontitle}
	\item[1.11]<1-> \hyperlink{1_11_a}{\IkSubsectiontitle}
\end{itemize}

%\endofslide 
\end{frame}
%%%%%%%%%%%%%%%

\xdef\sectiontitle{\IaSubsectiontitle} 

%%%%%%%%%%%%%%%
\begin{frame}%[label \equal 1_1_a]
\frametitle{\texorpdfstring{\culine[\sectiontitleulinecolor]{\sectiontitle}}{\sectiontitle} }

\begin{itemize}
	\item Statistical mechanics makes predictions about properties and behaviour of \appear{\Alert[red]{systems whose number of particles is enormous}}
	
\item The more particles\appear{, the more details are unknown}
\item For precise predictions\appear{, large number of particles improves the representativeness of averages}\appear{, and their evolution}
% \implies{ the more reliable are our predictions on evolution of the averages,  }
% \hspace{1cm}\implies{ the more precise are our predictions }

\item \CDAlert[red]{Thermodynamic system} is a system\appear{, whose number of particles is large enough for the predictions to become very precise}

\item Statistical mechanics:
	\item[] \begin{itemize}[itemsep=2.5mm]
		\item classical
		\item not 'modern' physics \appear{(i.\,e. not based on quantum mechanics)}
		\item BUT \CDAlert[tunired]{supports many areas of physics!}
	\end{itemize}
\end{itemize}

\endofslide 
%\endofsection
\end{frame}
%%%%%%%%%%%%%%%

%%%%%%%%%%%%%%%%
%\begin{frame}
%\frametitle{\texorpdfstring{\culine[\sectiontitleulinecolor]{\sectiontitle}}{\sectiontitle} }
%
%%\begin{minipage}{10.4cm}
%\begin{itemize}
%	\item Examine a two-sided room \appear{with $N$ point-like\\
%		 particles placed in the room} \appear{(see Fig.~1)}
%
%	\item For each particle\appear{, there are two alternatives:}\\ 
%		\appear{\Alert[red]{left}} \appear{or \Alert[blue]{right} } 
%	\item This is a thermodynamic \CDAlert[fuchsia]{'two-state system'} \\
%		\appear{(Recall coins or paramagnets in Thermal physics)}
%	\item Estimate in \CDAlert[fuchsia]{how many ways} we can place $N$ particles into the room \appear{so that $N_{R}$ number of particles are on the right-hand side?}
%	\item The answer is \CDAlert[fuchsia]{binomial coefficients}:	
%	\[
%\Omega \textrm{((} N, N_{R} \textrm{))}  \equal \appear{ W_{N_{R}}^{N} } \equal  \appear{\LP} 
%\begin{array}{l}
%N \\
%N_{R}
%\end{array} 
%\appear{\RP}   \equal \frac{N !}{N_{R} ! \textrm{((} N-N_{R} \textrm{))} !}
%\]
%\appear{giving us the '\Alert[tunired]{multiplicity}' or the ''\CDAlert[tunired]{number of ways to arrange...}"}
%
%\end{itemize}
%%\end{minipage}
%
%\xdef\loc{north east}
%\begin{tikzpicture}[remember picture,overlay] % anchor \equal south
%    \node (A) [yshift = 1*\figYshift,xshift = \figXshift, anchor = \loc] at (current page.\loc) {\includegraphics[page = 1,width = 4cm]{Figures_booklet/SOM_9_Statistical_mechanics_figures/SOM_Figure_1.png}}; %scale \equal \figscale. % 8cm max hei
%\end{tikzpicture}
%
%
%\endofslide 
%%\endofsection
%\end{frame}
%%%%%%%%%%%%%%%%


%%%%%%%%%%%%%%%
\begin{frame}
\frametitle{\texorpdfstring{\culine[\sectiontitleulinecolor]{\sectiontitle}}{\sectiontitle} }

%\begin{minipage}{10.4cm}
\begin{itemize}
	\item The ways how four \appear{(\CDAlert[fuchsia]{distinguishable})} \appear{particles can be arranged into a two-sided room:}
\appear{	\xdef\loc{center}
\begin{tikzpicture}[remember picture,overlay] % anchor \equal south
    \node (A) [yshift = 1*\figYshift,xshift = -3.5cm, anchor = \loc] at (current page.\loc) {\includegraphics[page = 1,width = 5cm]{Figures_booklet/SOM_9_Statistical_mechanics_figures/SOM_Figure_2_a.png}}; %scale \equal \figscale. % 8cm max hei
\end{tikzpicture}
}
\appear{%\xdef\loc{north east}
\begin{tikzpicture}[remember picture,overlay] % anchor \equal south
    \node (A) [yshift = 4*\figXshift,xshift = 3.5cm, anchor = \loc] at (current page.\loc) {\includegraphics[page = 1,width = 5cm]{Figures_booklet/SOM_9_Statistical_mechanics_figures/SOM_Figure_2_b.png}}; %scale \equal \figscale. % 8cm max hei
\end{tikzpicture}}

\end{itemize}
%\end{minipage}



\endofslide 
%\endofsection
\end{frame}
%%%%%%%%%%%%%%%


%%%%%%%%%%%%%%%
\begin{frame}
\frametitle{\texorpdfstring{\culine[\sectiontitleulinecolor]{\sectiontitle}}{\sectiontitle} }

\begin{minipage}{6.5cm}
\begin{itemize}[itemsep=1.9mm]
	\item In the case $N = 4$\appear{, there are altogether 5 alternative values for $N_{R}\;\appear{(\,= 0,1,2,3,4)$}}
	\item Distribution $\Omega = \Omega \textrm{((} N_{R} \textrm{))}$ shows a slight maximum at $\appear{N_{R}} \equal \appear{2} \equal \appear{1 / 2 N}$ \appear{(Fig.~3)}
	
	\item The distribution in the case $N = 4$ is broad\appear{, but if we increase value of $N$}\appear{, the \Alert[red]{distribution becomes narrower}}
	\end{itemize}
\end{minipage}

\vspace{1.9mm}
\begin{itemize}[itemsep=1.9mm]
	\item In all the cases, the maximum lies at $N_{R} = \frac{1}{2} \appear{N}$
	
	\item Note that for large equally-sized rooms and value of $N$\appear{, probably both rooms have almost exactly equal amount of particles}\appear{, although it would be impossible to know anything of a single molecule}
\end{itemize}


\xdef\loc{north east}
\begin{tikzpicture}[remember picture,overlay] % anchor \equal south
    \node (A) [yshift = 1*\figYshift,xshift = \figXshift, anchor = \loc] at (current page.\loc) {\includegraphics[page = 1,width = 7cm]{Figures_booklet/SOM_9_Statistical_mechanics_figures/SOM_Figure_3.png}}; %scale \equal \figscale. % 8cm max hei
\end{tikzpicture}

\endofslide 
%\endofsection
\end{frame}
%%%%%%%%%%%%%%%


%%%%%%%%%%%%%%%
\begin{frame}
\frametitle{\texorpdfstring{\culine[\sectiontitleulinecolor]{\sectiontitle}}{\sectiontitle} }

\begin{itemize}
	\item If we assume that the particles are air molecules in a room\appear{, and the two sides of the room are of equal size,} \appear{ then at highest probability \Alert[red]{both sides have almost exactly equal numbers of particles}} 
	\item This can be said with high relative accuracy/confidence...
	\vspace{4.5cm}
	\item[] ...although it would be impossible to know anything about the position of a single molecule
\end{itemize}

\xdef\loc{center}
\begin{tikzpicture}[remember picture,overlay] % anchor \equal south
    \node (A) [yshift = 1.45*\figYshift,xshift = 0cm, anchor = \loc] at (current page.\loc) {\includegraphics[page =1,width =12cm]{Figures_booklet/SOM_9_Statistical_mechanics_figures/FYS_1627_SSP-4.pdf}}; %scale \equal \figscale. % 8cm max hei
\end{tikzpicture}
\endofslide 
%\endofsection
\end{frame}
%%%%%%%%%%%%%%%

%%%%%%%%%%%%%%%
\begin{frame}
\frametitle{\texorpdfstring{\culine[\sectiontitleulinecolor]{\sectiontitle}}{\sectiontitle} }


%\begin{itemize}
	% 6.5 cm  \equal  midway, 10 cm  \equal  3/4 
%\item[] 
\begin{minipage}{8.5cm}
	\begin{itemize}[itemsep=1.8mm]
	\item The example above can be used to clarify
		 the terms \CDAlert[red]{microstate} \appear{and \CDAlert[blue]{macrostate}}

\item In Fig. 3d \appear{the highest peak of the distribution rises to value of $10^{3 \times 10^{22}}$} \appear{so, there are $10^{3 \times 10^{22}}$ different 'ways' of how
	 to arrange molecules in the room} \appear{if they are split into two equal}\\
	 \end{itemize}
\end{minipage}
\vspace{-4mm}
\begin{itemize}[itemsep=1.8mm]
	\item[]  fractions between left and right \textrm{((}specified by $ N_{\mathrm{R}} \textrm{))}$ 

\item Each 'way' is a different \CDAlert[red]{microstate}.
\appear{If we would know the states of all individual particles}\appear{, we could define each (micro)state}

\item All those $10^{3 \times 10^{22}}$ states \appear{belong into group where all the molecules are split into to equal halves.} \appear{Thus, all those $10^{3 \times 10^{22}}$ states are \CDAlert[blue]{microstates}, associated with $N_{R} \equal \frac{1}{2} \appear{N}$}

\item This is one \CDAlert[blue]{macrostate}. \appear{Others are those with different values of $N_{R}$}
\end{itemize}

\xdef\loc{north east}
\begin{tikzpicture}[remember picture,overlay] % anchor \equal south
    \node (A) [yshift =1*\figYshift,xshift=\figXshift, anchor=\loc] at (current page.\loc) {\includegraphics[page=1,width=4.8cm]{Figures_booklet/SOM_9_Statistical_mechanics_figures/SOM_Figure_3_d.png}}; %scale \equal \figscale. % 8cm max hei
\end{tikzpicture}

\endofslide 
%\endofsection
\end{frame}
%%%%%%%%%%%%%%%

%%%%%%%%%%%%%%%
\begin{frame}
\frametitle{\texorpdfstring{\culine[\sectiontitleulinecolor]{\sectiontitle}}{\sectiontitle} }
\seminormal

\begin{itemize}%[itemsep=2mm]
	\item Starting with the two-sided room and significantly different amounts of molecules on each side\appear{, the system \CDAlert[tunired]{evolves towards the case} \CDAlert[tunired]{of highest probability}}\appear{, i.\,e. towards the macrostate with highest number of 'ways'}\appear{, i.e. the one with molecules equally on both sides}

\item The most probable macrostate is known as the \CDAlert[red]{equilibrium state}

\item If left alone, a thermodynamic system will evolve towards equilibrium \appear{and eventually reach an equilibrium}
\end{itemize}

% 6.5 cm  \equal  midway, 10 cm  \equal  3/4 
\vspace{2.5mm}
\begin{minipage}{8.5cm}
\begin{itemize}%[itemsep=2mm]
	\item Starting with non-equilibrium\appear{, both macrostate and microstate of the system will change}

\item \CDAlert[blue]{At equilibrium}, the system still continues \Alert[tuniblue]{changing microstates}\appear{, but they all correspond to the \Alert[blue]{same macrostate}} \appear{(or the ones close to it...)}\appear{, macrostate specified by $ N_{\mathrm{R}}$}
\end{itemize}
\end{minipage}
\vspace{3mm}

\xdef\loc{south east}
\begin{tikzpicture}[remember picture,overlay] % anchor \equal south
    \node (A) [yshift=-0.25*\figYshift,xshift =\figXshift, anchor=\loc] at (current page.\loc) {\includegraphics[page=1,width=4.8cm]{Figures_booklet/SOM_9_Statistical_mechanics_figures/SOM_Figure_3_d.png}}; %scale \equal \figscale. % 8cm max hei
\end{tikzpicture}

%\endofslide 
\endofsection
\end{frame}
%%%%%%%%%%%%%%%



